\section*{Chapter 8: Gravitational Waves: Curvature In Motion}
\begin{description}[style=multiline,leftmargin=2em,labelwidth=1.5em]

% DOCX paragraph 2086
\item[1.]\hypertarget{note:08-01}{}A wave equation for a scalar field \(\phi\) can be written simply as
\(\ \square\ \phi\  = \ 0\), where \(\square\) is the \emph{D'Alambert differential
operator}\index{D’Alambert operator} (or simply D'alambertian) and is a synthetic representation
of a series of partial derivatives of the second order. For example, in a Cartesian
coordinate system in three spatial dimensions
\(\square\  = \partial_{t}^{2} - v^{2}(\partial_{x}^{2} + \partial_{y}^{2} + \partial_{z}^{2})\),
where \(v\) is the speed of propagation of the wave.

% DOCX paragraph 2087
\item[2.]\hypertarget{note:08-02}{}As mentioned in the text, the description I give of the motion of the mattress is valid
at the first order in the perturbation. If we wanted to also include nonlinear effects,
that is, contributions to the dynamics of the mattress of order higher than the first
one, the mattress would also move in the direction of propagation of the wave, as the
latter is partially longitudinal and so transfers a certain amount of linear momentum.
In practice, we would also drift slightly away from the person who jumped into the
water.

% DOCX paragraph 2088
\item[3.]\hypertarget{note:08-03}{}Alongside these two linear polarisations, there are also two \emph{circular
ones}\index{polarization!circular} (one for each direction of rotation) in which the deformation
produced by the gravitational wave rotates around the direction of propagation of the
wave. Referring again to Figure \ref{fig:08-02}, in this case, the Vitruvian\index{Leonardo da Vinci!Vitruvian man} Man
would be stretched and compressed along a direction that rotates over time as the
gravitational wave propagates.

% DOCX paragraph 2089
\item[4.]\hypertarget{note:08-04}{}In the very first moments of its life, the universe was extremely hot and dense, so that
matter and radiation were ``coupled'', and their temperatures evolved in the same way.
When the universe expanded and cooled enough, matter and radiation ``decoupled'' and, in
terms of thermal evolution, each ``went its own way.'' The moment of decoupling was very
special as the photons were finally ``free'' to propagate without constantly scattering
off all the particles around. These photons, which were only produced 380,000 years
after the Big Bang, are still visible in what is called the \emph{cosmic
microwave}\index{cosmic microwave background} \emph{background}, a very uniform bath of radiation. If we
had microwave-sensitive eyes, the night sky would not appear dark but dimly lit by this
primordial radiation.

% DOCX paragraph 2090
\item[5.]\hypertarget{note:08-05}{}While it is easy to define an event horizon in the case of a static or stationary
spacetime, it is much more difficult to calculate it in the case of a dynamic spacetime,
such as the one that characterises two black holes soon before merger, or a star before
collapsing into a black hole. In these cases, it is much more convenient to adopt the
concept of \emph{``apparent horizon''} to define the two-dimensional surface within
which the light is trapped at any given time. In the case of a static or stationary
spacetime, the two types of horizon coincide.

% DOCX paragraph 2091
\item[6.]\hypertarget{note:08-06}{}The Love number has nothing to do with feelings: it is named after the British physicist
Augustus Love\index{Love!Augustus} (1863-1940), who first defined it in his theory of
elasticity.

% DOCX paragraph 2092
\item[7.]\hypertarget{note:08-07}{}It goes without saying that if the mass of the object is above about 3 solar masses, it
will not be a neutron star but rather a black hole. However, there are observations
where the distinction is not so sharp. For example, in the event
GW190814\index{GW190814}, the least massive object has a mass of about 2.6 solar
masses. It is, therefore, not clear whether the observations refer to the least massive
black hole ever observed or to the most massive neutron star ever detected (see The LIGO
Collaboration and Virgo Collaboration, in \emph{Astrophysical Journal Letters}, 896, L44
(2020)).

% DOCX paragraph 2093
\item[8.]\hypertarget{note:08-08}{}Much of what we have seen for binary systems of black holes or neutron stars also
applies to a mixed binary system, that is, consisting of a black hole and a neutron
star. However, there are also two important differences. The first is related to the
mass ratio between the two objects, which can be 10 to 1 or even greater given the
characteristic masses involved (the black hole is expected to have a mass between 15-25
solar masses, while the neutron star between 1.3-2.3 solar masses). The second
difference is that the coalescence of a mixed system cannot lead to the formation of a
hypermassive star. In this case, the neutron star is either tidally destroyed and builds
up an accretion torus or is even ``absorbed whole'' if the difference in mass is very
large.

% DOCX paragraph 2094
\item[9.]\hypertarget{note:08-09}{}See E. Nakar, in \emph{Physics Reports}, 442, 166 (2007) for a review of the literature
on this subject.

% DOCX paragraph 2095
\item[10.]\hypertarget{note:08-10}{}See S. E. Woosley and J. S. Bloom, in \emph{Annual Review of Astronomy and
Astrophysics}, 44, 507 (2006) for a review of the relevant literature.

% DOCX paragraph 2096
\item[11.]\hypertarget{note:08-11}{}Gravitational-wave detectors are generically ``\emph{omnidirectional''}, that is,
sensitive in all directions, but not with the same sensitivity. In fact, there are
directions in which the sensitivity is maximum and others in which it is very small,
representing the detector's ``blind spots''. Hence, as with a radio antenna, it is
possible to define a sensitivity map of an interferometric gravitational-wave detector.
In the case of GW170817\index{GW170817}, the fact that the Virgo detector did not
measure the signal was nevertheless very valuable. The knowledge of Virgo's blind spot
at that time made it possible to narrow the search area for the electromagnetic signal
and, therefore, to discover GRB 170817A\index{GRB 170817A}.

% DOCX paragraph 2097
\item[12.]\hypertarget{note:08-12}{}In general, a strong electromagnetic counterpart\index{electromagnetic counterpart} is not expected from
the coalescence of binary systems of stellar-mass black holes. The reason is that these
systems represent the final stage of a binary system that has long since lost all traces
of matter, either through the collapse processes that produced the black holes or
through the supernova explosions and associated shock\index{shock} waves that
dispersed the matter in the interstellar space.

% DOCX paragraph 2098
\item[13.]\hypertarget{note:08-13}{}L. Rezzolla, B. Giacomazzo, L. Baiotti, J. Granot, K. Koveliotou, M.-A. Aloy, in
\emph{The Astrophysical Journal Letters}, 732, L6 (2011).

% DOCX paragraph 2099
\item[14.]\hypertarget{note:08-14}{}L.-X. Li and B. Paczyński, in \emph{The Astrophysical Journal Letters}, 507, L59 (1998).

% DOCX paragraph 2100
\item[15.]\hypertarget{note:08-15}{}In astronomy, it is referred to as \emph{nova}\index{nova}, a thermonuclear
explosion on the surface of a white dwarf, caused by the accumulation of hydrogen due to
an accretion process from a companion star in a binary system. The explosion can be
visible for several days and reaches a brightness that varies between 10,000 and 100,000
times that of the Sun. These explosions release minute amounts of energy when compared
to gamma-ray bursts.

% DOCX paragraph 2101
\item[16.]\hypertarget{note:08-16}{}J. M. Lattimer and D. N. Schramm, in \emph{The Astrophysical Journal}, 210, 549 (1976).

\end{description}
